\documentclass[linenumbers,amsmath,trackchanges]{aastex702}

\newcommand{\vdag}{(v)^\dagger}
\newcommand\aastex{AAS\TeX}
\newcommand\latex{La\TeX}

\begin{document}
\title{Placeholder}

% \author[orcid=0000-0000-0000-0001,sname='North America']{Tundra North America}
% \altaffiliation{Kitt Peak National Observatory}
% \affiliation{University of Saskatchewan}
% \email[show]{fakeemail1@google.com}  

% \author[gname=Savannah,sname=Africa]{S. Africa}
% \affiliation{South African Astronomical Observatory}
% \affiliation{University of Cape Town, Department of Astronomy}
% \email{fakeemail3@google.com}

% \collaboration{all}{The Terra Mater collaboration}

\begin{abstract}
What's the study about, why's it important, how to do it, what did you find, and why are the findings 
important?\\
\end{abstract}

%% https://astrothesaurus.org
% \keywords{\uat{Galaxies}{573} --- \uat{Cosmology}{343} --- \uat{High Energy astrophysics}{739} --- 
% \uat{Interstellar medium}{847} --- \uat{Stellar astronomy}{1583} --- \uat{Solar physics}{1476}}

To-do list:
\begin{itemize}
    \item Paper writing
    \begin{itemize}
        \item Short document with the outline of the paper and the main points to be included in each 
        section for people's ease of reading
        \item Methodology section: people should be able to replicate the results with the information 
        provided in this section; be concise; can direct readers to other papers for details
        \item Results section: analyse trends if there are any
        \item Why are the findings important: read Ankur's paper to see what people in this field care about
    \end{itemize}
    \item Figures
    \begin{itemize}
        \item Fig 1: recompute horizon redshift curve by setting all events to be face-on; cut high-mass 
        tail as snr is not computed properly which caused biases; check high horizon redshift events to 
        see some lie outside the horizon redshift curve (inclination?); look up FINDCHIRP paper for the 
        analytical formula for horizon distance
        \item Fig 2: delete?
        \item Fig 3: delete?
        \item Fig 4: abnormal redshift distribution (negative); change the axis from log mass to mass, 
        but change the scale to logarithmic to make things more intuitive; may need to use previous 
        version of population density for plotting this (\text{population\_densities\_3d.npy}); see if 
        can combine the detected and intrinsic populations to get rid of some white spaces
        \item Figs 5, 6, 7: check if the systems used for reconstructing the n-image lensed population 
        are truly $n$-image systems, or if some of them are $m$-image systems where $m>n$ and only $n$ 
        images are detectable; try to show dLs instead of ratios (same for time delays; higher dimension); 
        point out and try to explain correlations if they still exist after consistency checks
    \end{itemize}
    \item Consistency checks
    \begin{itemize}
        \item Models used by LeR to calculate rates and by Ankur to simulate GW signals
        \item Hypotheses: intrinsic n-image systems or systems having n images detectable; check Otto's
        prior odds paper, Ankur's PO2 script; can do both and they should give the same conclusion
        \item Check luminosity distance prior used in PO2: Madau and Dickinson for unlensed, lensed 
        population for lensed (again, check intrinsic or detectable)
    \end{itemize}
    \item Comparison: using uniform prior vs using reconstructed population as prior in PO2 to illustrate 
    the importance of population reconstruction (check if Ankur has it already - if so, no need to run 
    script again)
\end{itemize}


\section{Introduction}

\section{Methodology}
\subsection{Simulated GW events and Bayesian parameter estimation}
Do I have to elaborate on how we simulated the GW events, or briefly mention it and focus on how we use it?

\subsection{Population reconstruction}
\label{sec:population_reconstruction}
% Natalie's part
The population distribution of binary black hole (BBH) parameters plays a crucial role in determining 
whether two similar gravitational wave (GW) signals originate from the same source or from two independent 
mergers exhibiting very similar properties.\\
In densely populated regions of the parameter space, two unrelated GW events are more likely to exhibit 
similar inferred parameters purely by coincidence. Conversely, agreement between parameters that are rare 
in the underlying BBH population provides stronger evidence in favour of the lensing hypothesis.\\
To account for this population dependence in the Bayesian analysis, we reconstruct the BBH population 
distribution from simulated GW observations. We first reconstruct the observed population using posterior 
samples from detected events, then correct for selection effects to recover the underlying astrophysical 
population. The reconstructed population is subsequently used to generate a strongly lensed BBH population 
for prior odds estimation and Bayes factor computation.

\subsubsection{Reconstruction of observed population}
Population reconstruction on a simulated distribution of BBHs is performed using FIGARO, which uses a 
non-parametric approach to model probability densities using a Dirichlet Process Gaussian Mixture Model 
(DPGMM). The probability density is represented as a mixture of Gaussian components, with the number of 
components inferred from the data rather than specified in advance. Since the parameters of individual 
GW events are inferred with measurement uncertainties, we employ the hierarchical DPGMM (HDPGMM) framework 
to combine information from individual-event posterior distributions and reconstruct the underlying 
population density~\citep{RinaldiDelPozzo2022}.\\
The reconstruction is performed on the joint three-dimensional distribution of the component masses $m_1$, 
$m_2$, and redshift $z$. For each event, the posterior density is first reconstructed using DPGMM, with 
multiple density realizations generated to characterize the uncertainty in the density estimate. These 
individual-event density reconstructions are subsequently combined hierarchically using HDPGMM to infer 
the population distribution
(see Fig.~\ref{fig:reconstructed_population}, left panel, for the reconstructed population). 

\subsubsection{Correction for selection effects}
The reconstructed observed population does not reflect the intrinsic population, as the reconstruction is 
based on posterior samples of detected events. In order to eliminate selection effects and create an 
astrophysical population, we perform rejection sampling during generation of samples from the reconstructed 
observational population, with acceptance probability approximated to be a function of chirp mass.\\
The astrophysical BBH merger rate density is assumed to follow a redshift-dependent model,
\[\frac{R_\text{GW}(z)}{\text{Gpc}^{-3}\text{yr}^{-1}} = \frac{a_1 e^{a_2 z}}{e^{a_3 z}+a_4},\]
where $a_1=6600$, $a_2=1.6$, $a_3=2.1$, and $a_4=30$ for population-I/II stars\citep{Oguri2018}.\\
The number of detectable events for a given chirp mass is given by
\[N_\text{detected}(M_c) = \int_0^{z_\text{horizon}(M_c)}\frac{dV_c}{dz}\frac{R(z)}{1+z}dz,\]
where $\frac{dV_c}{dz}$ is the differential comoving volume element and the factor (1/(1+z)) accounts 
for cosmological time dilation.\\
In importance sampling, the importance ratio, or the weight, is given by 
\[w(x) \propto \frac{p(x)}{q(x)}, \ x_i\sim q(x),\]
where $p(x)$ and $q(x)$ are the target population and biased population respectively.\\
In our case, it can written as
\[w(M_c) \propto \frac{P_\text{intrinsic}(M_c)}{P_\text{FIGARO}(M_c)}, \ M_c\sim P_\text{FIGARO}(M_c),\]
which is equivalent to
\[w(M_c) \propto \frac{P(M_c)}{P(M_c | \text{detected})}.\]
Applying Bayes theorem, 
\[P(M_c | \text{detected}) = \frac{P(\text{detected} | M_c) P(M_c)}{P(\text{detected})},\]
where $P(\text{detected})$, the probability of detection, is independent of chirp mass, 
and \[P(\text{detected} | M_c) = \frac{N_\text{detected}(M_c)}{N_\text{all}}.\]
Hence, the reweighting factor can be written as
\[w(M_c) \propto \frac{1}{N_\text{detected}(M_c)},\]
which is then converted to acceptance probability $a(M_c)$ afterwards for rejection sampling:
\[a(M_c) = \frac{w(M_c)}{w_\text{max}}.\]
Rejection sampling is then performed on draws from the observed population to obtain an approximately 
selection-corrected astrophysical population, which is subsequently used to compute the prior odds and 
Bayes factors.\\
The correction generally favours lower-mass systems, which are detectable over smaller cosmological volumes. 
This trend is important to capture when imposing strong lensing on the reconstructed population (see
Fig.~\ref{fig:reconstructed_population}, right panel, for the reconstructed
population).

\subsubsection{Reconstruction of the strongly lensed population}
% Natalie: To fill: 
% - How we lens the population
% - two-image, three-image, four-image cases
% - Figure: (effective luminosity distances and log arrival times) + caption
To reconstruct the strongly lensed population, we use the LVK (LIGO–Virgo–KAGRA) Event Rate calculator and 
simulator (LeR) \cite{}.
Component masses are drawn from the reconstructed BBH distribution corrected for selection effects, 
while source redshifts are sampled according to the astrophysical merger-rate distribution. 
The remaining GW source parameters are sampled from their respective distributions. For each source, 
a foreground lens galaxy is generated using LeR, with lens redshift, velocity dispersion, ellipticity, 
orientation, and external shear sampled according to the lens population model. The lens equation is then 
solved to obtain the image positions, magnifications, and arrival-time delays. The resulting lensing 
parameters are stored for subsequent analysis, and detectable lensed events are selected using LeR's GW 
detection-probability calculation.\\
The number of lensed images depends on the source position relative to the caustics, which are 
boundaries in the source plane separating regions with different image multiplicities. Sources 
located between the outer and inner caustics generally produce two images, while those inside the 
inner diamond-shaped caustic produce four images. For sufficiently elliptical lenses, the diamond 
caustic may extend beyond the outer caustic, allowing three-image configurations. To allow further 
investigation and comparison between different cases, we reconstruct the population for 2-image, 
3-image, and 4-image cases separately. For each image multiplicity, we reconstruct the joint distribution 
of relative effective luminosity distances and arrival-time delays using FIGARO's DPGMM
(see Fig.~\ref{fig:lensed_population_l2},~\ref{fig:lensed_population_l3}, and~\ref{fig:lensed_population_l4} 
for the reconstructed lensed population).

\subsection{Prior odds estimation}
\label{sec:prior_odds}
% Natalie's part
The prior odds are defined as the ratio of the probability that a given set of GW events consists entirely 
of strongly lensed images of the same source to the probability that this is not the case, rather than the 
ratio of the probability that an individual event is lensed to the probability that it is not~\citep{Hannuksela2026}. 
As we have already corrected for selection effects in the population reconstruction, the prior odds are 
estimated using the detectable merger rates obtained from the reweighted population:
\[
P_U^L = \frac{P(H_L)}{P(H_U)} \approx P(H_L) \approx N_{\text{img}}! 
\frac{\langle N_L \rangle}{\langle N \rangle ^{N_{\text{img}}}},
\]
where $N_{\text{img}}$ is the number of images, $\langle N_L \rangle = R_L T$ is the expected number 
of detectable lensed systems, and $\langle N \rangle = (R_U+N_{\text{img}}R_L)T$ is the expected total 
number of detected GW signals over an observation time $T$. Here, $R_L$ and $R_U$ denote the detectable 
lensed-system and unlensed-event rates, respectively~\citep{Hannuksela2026}.\\
The detectable rates are estimated using LeR from the simulated unlensed and lensed populations 
described in Section~\ref{sec:population_reconstruction}. To distinguish different image multiplicities, 
the detectable lensed systems are classified into 2-image, 3-image, and 4-image configurations. For 
each configuration, we require that all generated images satisfy the detection criterion. The rate 
corresponding to each configuration is estimated by weighting the total detectable lensed rate by 
the fraction of systems satisfying both the image-multiplicity and detectability requirements,
\[
R_{L,n} = R_L\frac{N_{L,n}^{\text{all\,det}}}{N_L^{\text{det}}},
\]
and systems for which only a subset of the generated images is detectable are excluded from the 
multiplicity-specific rates.\\
The prior odds for each image multiplicity are then calculated as
\[
P_{U,n}^{L} \approx n!\frac{R_{L,n}T}{\left[(R_U+nR_{L,n})T\right]^n}, \qquad n=2,3,4.
\]
Since the lensed-event rate is much smaller than the unlensed-event rate, \(R_{L,n}\ll R_U\), this can 
be approximated as
\[
P_{U,n}^{L} \approx n!\frac{R_{L,n}}{R_U^n T^{n-1}}.
\]
We adopt an observation time of \(T=1,\mathrm{yr}\) when evaluating the prior odds for the 2-, 3-, and 
4-image cases.

\subsection{Bayes factor computation}
% Natalie's part
We adopt the PO2.0 framework developed by \citet{Barsode2025} to compute the Bayes factor between 
the strongly lensed and unlensed hypotheses. The Bayes factor is defined as
\[
B_U^L = \frac{P(d_1,\ldots,d_N\mid H_L)}{P(d_1,\ldots,d_N\mid H_U)},
\]
where \(d_i\) denotes the GW data associated with the \(i\)-th event and \(N\) is the number of 
candidate images. PO2.0 evaluates the Bayes factor using posterior samples from individual-event 
parameter estimation, without requiring joint parameter estimation. The calculation incorporates 
the consistency of the intrinsic binary parameters and sky localization, together with lensing-induced 
differences in apparent luminosity distance, arrival time, and phase. We refer readers to 
\citet{Barsode2025} for the detailed formulation.\\
As suggested by \citet{Barsode2025}, we generalize the original two-event PO2.0 implementation to 
accommodate \(N\)-image systems. The events are ordered by arrival time, with the earliest-arriving 
image taken as the reference. The posterior samples of the reference image are transformed according to 
the relative magnifications and Morse phase differences, and compared with the posterior distributions 
of the remaining \(N-1\) images. The calculation is extended to include the joint distributions of 
the relative lensing parameters, and contributions from different Morse phase configurations are 
combined according to their relative probabilities in the simulated lensed population.\\
Following Equation B15 of \citet{Barsode2025}, we generalize the time-delay likelihood ratio to 
\(N\) candidate images. Under the unlensed hypothesis, the arrival times are assumed to be independently 
and uniformly distributed within an observation period \(T_{\mathrm{obs}}\). For \(N\) events ordered 
by arrival time, \(t_1<t_2<\cdots<t_N\), the joint probability density is
\[
P(t_1,\ldots,t_N\mid H_U)=\frac{N!}{T_{\mathrm{obs}}^N}.
\]
Defining the relative time delays as \(\Delta t_j=t_j-t_1\), where \(j=2,\ldots,N\), and marginalizing 
over the first arrival time \(0<t_1<T_{\mathrm{obs}}-\Delta t_N\), we obtain
\[
P(\Delta t_2,\ldots,\Delta t_N\mid H_U) = \frac{N!(T_{\mathrm{obs}}-\Delta t_N)}{T_{\mathrm{obs}}^N}.
\]
Since the reconstruction is performed using logarithmic time delays, we define the vector 
\(\boldsymbol{x}=(x_2,\ldots,x_N)\), where \(x_j=\log_{10}(\Delta t_j/\mathrm{s})\). Applying the 
corresponding Jacobian, we obtain
\[
P(\boldsymbol{x}\mid H_U) = \frac{N!(T_{\mathrm{obs}}-\Delta t_N)}{T_{\mathrm{obs}}^N}
(\ln 10)^{N-1}\prod_{j=2}^{N}\Delta t_j.
\]
The corresponding joint distribution under the lensed hypothesis is estimated using kernel density 
estimation from the simulated lensed population. The time-delay likelihood ratio is therefore
\[
\mathcal R_{LU}^{(N)} = \frac{P(\boldsymbol{x}\mid H_L)}{P(\boldsymbol{x}\mid H_U)}.
\]
We adopt \(T_{\mathrm{obs}}=1\,\mathrm{yr}\), consistent with the observation time used for the 
prior-odds calculation. 
% To account for the priors adopted during individual-event parameter estimation 
% with Cogwheel, we construct corresponding prior dictionaries using Bilby's BBHPriorDict. Uniform 
% detector-frame component-mass priors are imposed together with the chirp-mass and mass-ratio constraints, 
% while the luminosity-distance prior is taken to be proportional to \(d_L^2\). These reproduce the 
% principal functional dependencies of the Cogwheel PE priors, although exact equivalence has not been 
% independently established.\\
Finally, we automate the calculation over the simulated catalogues of 2-, 3-, and 4-image systems. For 
each system, the individual-event posterior samples and corresponding PE priors are loaded, preliminary 
posterior-overlap checks are performed, and the lensed and unlensed evidences are evaluated. For 
numerical stability, the logarithmic Bayes factors are restricted to the range $[-30,30]$, with systems 
failing the preliminary posterior-overlap checks assigned $\log_{10} B_U^L = -30$. The resulting Bayes 
factors are combined with the prior odds calculated in Section~\ref{sec:prior_odds} to obtain the 
posterior odds.

\section{Result}
\subsection{Reconstructed Population}
% Natalie's part
\begin{figure}[h!]
    \centering
    \includegraphics[width=0.9\textwidth]{figures/z_hor_vs_chirp_mass.pdf}
    \caption{Horizon redshift $z_{\mathrm{horizon}}$ as a function of chirp mass $\mathcal{M}_c$. The 
    scattered points represent the simulated detectable binary black hole events, while the solid curve 
    shows the estimated horizon redshift for each chirp mass. The horizon redshift increases with chirp 
    mass up to approximately $180\,M_\odot$, before decreasing at higher masses, reflecting the 
    mass-dependent sensitivity of gravitational-wave detectors.}
    \label{fig:z_hor_vs_chirp_mass}
\end{figure}
\begin{figure}[h!]
    \centering
    \includegraphics[width=0.9\textwidth]{figures/log_weight_vs_chirp_mass.pdf}
    \caption{Logarithm of the selection-correction weight as a function of chirp mass $\mathcal{M}_c$, 
    calculated from the expected number of detectable mergers within the horizon redshift. The weight 
    decreases sharply with increasing chirp mass at lower masses, reaching a minimum around 
    $180$--$200\,M_\odot$, before increasing gradually at higher masses. Larger weights correspond to 
    systems with lower detection efficiency and therefore a greater probability of being retained during 
    the rejection sampling used to reconstruct the astrophysical population.}
    \label{fig:log_weight_vs_chirp_mass}
\end{figure}
\begin{figure}[h!]
    \centering
    \includegraphics[width=0.45\textwidth]{figures/unlensed_2d.pdf}
    \includegraphics[width=0.463\textwidth]{figures/population_reweighted_m_1_m_2.pdf}
    \caption{
    	Reconstructed observed (left panel) and astrophysical (right panel) binary black hole population in component masses $m_1$ and $m_2$ using the hierarchical DPGMM implementation in FIGARO inferred from simulated population of gravitational-wave signals mimicking observing run O4. 
	The diagonal panels show the one-dimensional marginal distributions, while the off-diagonal panels show the corresponding two-dimensional joint distributions. 
	The observed population is biased towards higher masses due to selection effects, while the astrophysical population is corrected for these effects and represents the true underlying distribution of binary black hole mergers, extending to higher redshifts and lower masses that are relevant to strong gravitational-wave lensing, which occurs for binaries at high redshift.
	}
    \label{fig:reconstructed_population}
\end{figure}
%\begin{figure}[h!]
%    \centering
%    \includegraphics[width=0.9\textwidth]{figures/unlensed_3d.pdf}
%    \caption{}
%    \label{fig:unlensed_3d}
%\end{figure}
\begin{figure}[h!]
    \centering
    \includegraphics[width=0.45\textwidth]{figures/unlensed_3d_log.pdf}
    \includegraphics[width=0.463\textwidth]{figures/population_reweighted_m_1_m_2_log.pdf}
    \caption{
    	Reconstructed observed (left panel) and astrophysical (right panel) binary black hole population in log10 of component masses $m_1$ and $m_2$ using the hierarchical DPGMM implementation in FIGARO inferred from simulated population of gravitational-wave signals mimicking observing run O4. 
	The diagonal panels show the one-dimensional marginal distributions, while the off-diagonal panels show the corresponding two-dimensional joint distributions. 
	The observed population is biased towards higher masses due to selection effects, while the astrophysical population is corrected for these effects and represents the true underlying distribution of binary black hole mergers, extending to higher redshifts and lower masses that are relevant to strong gravitational-wave lensing, which occurs for binaries at high redshift.
	}
    \label{fig:reconstructed_population_log}
\end{figure}
\begin{figure}[h!]
    \centering
    \includegraphics[width=0.6\textwidth]{figures/lensed_population_l2_dL_dt.pdf}
    \caption{Reconstructed strongly lensed binary black hole population for 2-image systems using 
    FIGARO's DPGMM. The reconstruction is performed on the joint distribution of the logarithmic relative 
    effective luminosity distance $\log_{10}(d_{L,2}^{\mathrm{eff}}/d_{L,1}^{\mathrm{eff}})$ and 
    logarithmic arrival-time delay $\log_{10}(\Delta t_2/\mathrm{s})$, with the earliest-arriving
    image taken as the reference. The diagonal panels show the one-dimensional marginal distributions, 
    while the off-diagonal panel shows the corresponding two-dimensional joint distribution. The 
    reconstructed density reproduces the main features of the simulated lensing-parameter distribution.}
	\label{fig:lensed_population_l2}
\end{figure}
\begin{figure}[h!]
    \centering
    \includegraphics[width=0.9\textwidth]{figures/lensed_population_l3_dL_dt.pdf}
    \caption{Reconstructed strongly lensed binary black hole population for 3-image systems using 
    FIGARO's DPGMM. The reconstruction is performed on the joint four-dimensional distribution of 
    logarithmic relative effective luminosity distances $\log_{10}(d_{L,j}^{\mathrm{eff}}/d_{L,1}^
    {\mathrm{eff}})$ and logarithmic arrival-time delays $\log_{10}(\Delta t_j/\mathrm{s})$, where 
    $j=2,3$. The images are ordered by arrival time, with the earliest-arriving image taken as the 
    reference. The diagonal panels show the one-dimensional marginal distributions, while the 
    off-diagonal panels show the corresponding two-dimensional joint distributions. The reconstructed 
    density captures the main features and correlations of the simulated lensing parameters.}
    \label{fig:lensed_population_l3}
\end{figure}
\begin{figure}[h!]
    \centering
    \includegraphics[width=0.9\textwidth]{figures/lensed_population_l4_dL_dt.pdf}
    \caption{Reconstructed strongly lensed binary black hole population for 4-image systems using 
    FIGARO's DPGMM. The reconstruction is performed on the joint six-dimensional distribution of 
    logarithmic relative effective luminosity distances $\log_{10}(d_{L,j}^{\mathrm{eff}}/d_{L,1}^
    {\mathrm{eff}})$ and logarithmic arrival-time delays $\log_{10}(\Delta t_j/\mathrm{s})$, where 
    $j=2,3,4$. The images are ordered by arrival time, with the earliest-arriving image taken as the 
    reference. The diagonal panels show the one-dimensional marginal distributions, while the 
    off-diagonal panels show the corresponding two-dimensional joint distributions. The reconstructed 
    density captures the main features and correlations of the simulated lensing parameters.}
    \label{fig:lensed_population_l4}
\end{figure}

\subsection{Prior Odds}
% Natalie's part
%# Results:
%# Total unlensed rate: 448.76 detectable events per year
%# Total lensed rate: 0.44 detectable events per year
%# Rate ratio (lensed/unlensed): 0.00099
%# 2-image lensed rate: 0.19 detectable events per year
%# 2-image rate ratio: 0.00042
%# 2-image prior odds: 1.9e-06
%# 3-image lensed rate: 0.02 detectable events per year
%# 3-image rate ratio: 5e-05
%# 3-image prior odds: 1.5e-09
%# 4-image lensed rate: 0.07 detectable events per year
%# 4-image rate ratio: 0.00015
%# 4-image prior odds: 3.9e-11
\begin{table}[h!]
    \centering
    \caption{
    Estimated detectable gravitational-wave event rates,
    lensed-to-unlensed rate ratios, and prior odds for
    2-image, 3-image, and 4-image systems.
    The total detectable unlensed-event rate is
    $R_U=448.76\,\mathrm{yr}^{-1}$, while the total
    detectable lensed-system rate is
    $R_L=0.44\,\mathrm{yr}^{-1}$.
    The multiplicity-specific rates require all generated
    images to satisfy the GW detection criterion.
    The prior odds are evaluated for an observation time
    of $T=1\,\mathrm{yr}$.
    }
    \begin{tabular}{lccc}
        \hline
        Quantity & 2 images & 3 images & 4 images \\
        \hline
        $R_{L,n}\,(\mathrm{yr}^{-1})$
        & 0.19 & 0.02 & 0.07 \\

        $R_{L,n}/R_U$
        & $4.2\times10^{-4}$
        & $5.0\times10^{-5}$
        & $1.5\times10^{-4}$ \\

        $P_{U,n}^{L}$
        & $1.9\times10^{-6}$
        & $1.5\times10^{-9}$
        & $3.9\times10^{-11}$ \\
        \hline
    \end{tabular}
    \label{tab:prior_odds}
\end{table}


\subsection{Bayes Factor and Comparison}
% Natalie's part
\begin{figure}[h!]
    \centering
    \includegraphics[width=0.9\textwidth]{figures/bayes_factor_and_posterior_odds_comparison.pdf}
    \caption{Distributions of the logarithmic Bayes factors $\log_{10}B_U^L$ (upper panel) and posterior 
    odds $\log_{10}O_U^L$ (lower panel) for simulated strongly lensed 2-image (blue), 3-image (orange), 
    and 4-image (green) systems. The dashed vertical lines indicate the corresponding distribution peaks.
    The Bayes factor distributions shift towards larger values with increasing image multiplicity, 
    reflecting the additional information provided by multiple lensed images. After incorporating the 
    prior odds, the posterior odds distributions shift towards lower values due to the low astrophysical 
    probability of strong lensing.}
    \label{fig:log_bayes_factor_posterior_odds_gw}
\end{figure}
\begin{figure}[h!]
    \centering
    \includegraphics[width=0.9\textwidth]{figures/log_bayes_factor_vs_component_masses_all_images.pdf}
    \caption{Logarithmic Bayes factors $\log_{10}B_U^L$ as functions of component masses $m_1$ (left 
    panel) and $m_2$ (right panel) for simulated strongly lensed 2-image (blue), 3-image (orange), and 
    4-image (green) systems. The Bayes factors are generally larger for systems with higher image 
    multiplicities, with the highest values predominantly occurring at lower component masses. The 
    distributions show considerable scatter and overlap between image multiplicities. Values at 
    $\log_{10}B_U^L=-30$ correspond to the imposed numerical cutoff, including systems failing the 
    preliminary posterior-overlap checks.}
    \label{fig:scatter_log_bayes_factor_masses_gw}
\end{figure}


\section{Conclusion}











% \section{Software and third party data repository citations} \label{sec:cite}

% The AAS Journals would like to encourage authors to change software and
% third party data repository references from the current standard of a
% footnote to a first class citation in the bibliography.  As a bibliographic
% citation these important references will be more easily captured and credit
% will be given to the appropriate people.

% The first step to making this happen is to have the data or software in
% a long term repository that has made these items available via a persistent
% identifier like a Digital Object Identifier (DOI).  A list of repositories
% that satisfy this criteria plus each one's pros and cons are given at \break
% \url{https://github.com/AASJournals/Tutorials/tree/master/Repositories}.

% In the bibliography the format for data or code follows this format: \\

% \noindent author year, title, version, publisher, prefix:identifier\\

% \citet{2015ApJ...805...23C} provides a example of how the citation in the article references the external 
% code at \doi{10.5281/zenodo.15991}.  Unfortunately, bibtex does not have specific bibtex entries for these 
% types of references so the ``@misc'' type should be used. The Repository tutorial explains how to code the 
% ``@misc'' type correctly.  The most recent .bst file, aasjournalv7.bst, will output bibtex ``@misc'' type 
% properly.

% Authors can also use the website \url{https://www.doi2bib.org/} to create a BIBTeX entry for any DOI. 
% Please check the output from this site carefully as its output is only as good as the DOI metadata. 
% Some DOI creators do not provide enough metadata to construct an adequate citation.

%% Please use the acknowledgment and contribution environments. This will 
%% be anonomyized when the "anonymous" style option is used. 
\begin{acknowledgments}

\end{acknowledgments}

\begin{contribution}

\end{contribution}

% \facilities{HST(STIS), Swift(XRT and UVOT), AAVSO, CTIO:1.3m, CTIO:1.5m, CXO}

% \software{astropy \citep{2013A&A...558A..33A,2018AJ....156..123A,2022ApJ...935..167A},  
%           Cloudy \citep{2013RMxAA..49..137F}, 
%           Source Extractor \citep{1996A&AS..117..393B}
%           }

\appendix
\section{Placeholder appendix} \label{sec:appendix}

\bibliography{references}{}
\bibliographystyle{aasjournalv7.1}

%\allauthors

%\listofchanges

\end{document}


